\section{Threats to Validity}
\label{sec:threats}

\paragraph{Construct validity: KT-as-simulator}
Our sequencing-utility evaluation relies on a trained KT model acting as
a learner simulator. Such simulators are imperfect proxies for real
learners, and offline policy evaluation can be biased when the logged
policy differs from the evaluated one. We mitigate this with robustness
checks (multiple simulators, sensitivity analysis) and interpret
sequencing results as offline evidence rather than deployment guarantees.

\paragraph{Internal validity: leakage and splits}
KT evaluation is susceptible to label leakage and to student overlap
between train and test. We use student-level cross-validation and follow
established preprocessing to reduce these risks.

\paragraph{External validity: dataset and domain scope}
Findings are drawn from a limited set of public programming datasets
(primarily CS1 Java and intro Python) and one mathematics contrast; they
may not transfer to other languages, levels, or populations.

\paragraph{Reproducibility}
% \todo{Confirm released code, fixed seeds, and environment specification.}
We release code, configurations, and fixed random seeds; any AI
assistance and third-party models/datasets are declared with usage
evidence per the RBL submission guideline.
